%% ch02-installation.tex — JA4proxy Operator's User Guide
%% Chapter 2: Installation

\chapter{Installation}
\label{ch:installation}
\index{installation}

\newthought{JA4proxy runs as a Docker Compose stack.} There is no native package
installation and no binary release to configure by hand. Everything — the proxy
instances, Redis, HAProxy, the analytics node, and the observability stack — runs
in containers managed by Compose.

This chapter covers getting that stack running from a clean machine. It takes
approximately 10--15 minutes.

\section{Prerequisites}
\index{prerequisites}

\subsection{Software}

\begin{table}[ht]
  \small
  \begin{tabularx}{\linewidth}{lll}
    \toprule
    \textbf{Requirement} & \textbf{Minimum version} & \textbf{Notes} \\
    \midrule
    Docker Engine & 24.0 & Docker Desktop on macOS/Windows works \\
    Docker Compose & v2.20 & Bundled with recent Docker installations \\
    GNU Make & 3.81 & Pre-installed on Linux and macOS \\
    Git & Any & For cloning the repository \\
    \bottomrule
  \end{tabularx}
  \caption{Software prerequisites}
\end{table}

Verify your installed versions:

\begin{bashcode}
docker --version
docker compose version
make --version
\end{bashcode}

\subsection{Ports}

The following ports must be available on the host machine. If any are in use by
another service, stop that service or adjust the port mapping in
\code{docker-compose.yml} before proceeding.

\begin{table}[ht]
  \small
  \begin{tabularx}{\linewidth}{llX}
    \toprule
    \textbf{Port} & \textbf{Service} & \textbf{Purpose} \\
    \midrule
    443 & HAProxy & Public TLS entry point \\
    8080 & JA4proxy & Direct proxy access (internal) \\
    8404 & HAProxy stats & HAProxy statistics page \\
    8443 & Backend & Test backend HTTPS server \\
    8888 & Tarpit & Slow-response tarpit service \\
    9090 & JA4proxy metrics & Prometheus scrape endpoint (proxy) \\
    9091 & Prometheus & Prometheus query UI \\
    3001 & Grafana & Dashboard UI \\
    3100 & Loki & Log aggregation \\
    9093 & Alertmanager & Alert routing \\
    \bottomrule
  \end{tabularx}
  \caption{Port assignments}
\end{table}

\begin{notebox}[Firewall note]
In production, only port 443 should be exposed to the public internet. All other
ports should be restricted to your management network or accessible only via VPN.
The development stack exposes all ports to localhost for convenience.
\end{notebox}

\subsection{System Resources}

The full stack (proxy, Redis, analytics, observability) requires:
\begin{itemize}
  \item 4 GB RAM available to Docker
  \item 2 CPU cores
  \item 5 GB disk (images + GeoIP databases + logs)
\end{itemize}

On macOS and Windows, ensure Docker Desktop is configured with at least these
resources in its settings.

\section{Step-by-Step Installation}
\index{installation!steps}

\subsection{Step 1: Clone the Repository}

\begin{bashcode}
git clone https://github.com/seanpor/JA4proxy.git
cd JA4proxy
\end{bashcode}

\subsection{Step 2: Create the Environment File}
\index{.env file}

\ja\ uses a \code{.env} file for deployment-specific configuration that you do not
want to commit to version control (credentials, backend hostnames, API keys).

\begin{bashcode}
cp .env.example .env
\end{bashcode}

Open \code{.env} in your editor and set at minimum:

\begin{yamlcode}
# The hostname or IP address of your backend server
BACKEND_HOST=your.backend.host

# Grafana admin credentials
GRAFANA_ADMIN_USER=admin
GRAFANA_ADMIN_PASSWORD=changeme

# AbuseIPDB API key (optional — features degrade gracefully without it)
ABUSEIPDB_API_KEY=

# IP2Location API key for GeoIP database download
# Register free at https://lite.ip2location.com
IP2LOCATION_TOKEN=
\end{yamlcode}

\begin{warningbox}[Change default credentials]
The default Grafana credentials in \code{.env.example} are public. Change them
before running the stack, even in a development environment.
\end{warningbox}

\subsection{Step 3: Start the Stack}
\index{make start}

\begin{bashcode}
make start
\end{bashcode}

This builds all Docker images and starts the full stack. On a fast connection,
the first build takes 3--5 minutes as it pulls base images. Subsequent starts
are much faster.

You will see Docker Compose output as each container starts. When all containers
report healthy status, the stack is ready.

\subsection{Step 4: Download the GeoIP Database}
\index{GeoIP}
\index{IP2Location}
\index{make update-geoip}

Country-based filtering requires the IP2Location Lite country database. Download
and install it:

\begin{bashcode}
make update-geoip
\end{bashcode}

This downloads the database from IP2Location using the token in your \code{.env}
file and places it where the proxy can read it. If \code{IP2LOCATION\_TOKEN} is not
set, the command will print instructions for manual download.

\begin{notebox}[Free tier is sufficient]
The IP2Location Lite database is free for non-commercial use. Country-level
accuracy is what \ja\ needs — the free tier provides this. Register at
\url{https://lite.ip2location.com} and copy your token to \code{.env}.
\end{notebox}

Country filtering is disabled by default in \code{config/proxy.yml} until the
database is present. The proxy logs a warning at startup if GeoIP is enabled but
the database file is missing.

\section{Verifying the Installation}
\index{health check}
\index{verification}

\subsection{Container Status}

Check that all containers are running and healthy:

\begin{bashcode}
docker compose ps
\end{bashcode}

Expected output shows all services in the \code{running} state. The proxy and
Redis containers include health checks — wait for them to show \code{(healthy)}
before proceeding.

\subsection{Proxy Health Endpoint}

\begin{bashcode}
curl -s http://localhost:8080/health | python3 -m json.tool
\end{bashcode}

A healthy response:
\begin{yamlcode}
{
  "status": "healthy",
  "redis": "connected",
  "geoip": "loaded",
  "uptime_seconds": 42
}
\end{yamlcode}

\subsection{HAProxy Statistics}

Open \url{http://localhost:8404/stats} in a browser. You should see the HAProxy
statistics page showing the \code{ja4proxy\_backend} farm with at least one
server in UP state.

\subsection{Prometheus Scrape Targets}

Open \url{http://localhost:9091/targets}. All targets should show \code{UP}
state. If the \code{ja4proxy} target is down, the proxy metrics endpoint is
not responding — check the proxy container logs.

\subsection{Grafana Dashboard}

Open \url{http://localhost:3001} and log in with your Grafana credentials. The
\emph{JA4proxy Overview} dashboard should be available under Dashboards. At this
point it will show zero connections — that is expected.

\begin{tipbox}[Generate test traffic]
To see the dashboard populate immediately, generate some test traffic:
\begin{lstlisting}[language=bash, numbers=none]
./scripts/generate-tls-traffic.sh 30 50 5
\end{lstlisting}
This runs for 30 seconds with 50\% legitimate traffic and 5 worker threads.
See Chapter~\ref{ch:firststeps} for a full explanation of the traffic generator.
\end{tipbox}

\section{Rebuilding the Stack}
\index{make rebuild}

If you need to rebuild all images from scratch (after pulling a new version of the
repository, or if the images are in an inconsistent state):

\begin{bashcode}
make rebuild
\end{bashcode}

\begin{dangerbox}[Wipes all data]
\code{make rebuild} stops all containers, removes images, and rebuilds from scratch.
Any data in Redis (bans, rate-limit windows) will be lost. Use this only when you
need a clean slate or after updating the repository.
\end{dangerbox}

\section{Building the Go Proxy (Optional)}
\index{Go proxy}
\index{Phase 15}

A Go implementation of the proxy is included in \code{proxy\_go/}. It is currently
in active development targeting 10--50× the throughput of the Python implementation.
The Python proxy is the production-stable version; the Go proxy is for evaluation
and performance testing.

To build the Go proxy:

\begin{bashcode}
cd proxy_go
go build -o ja4proxy ./cmd/proxy
./ja4proxy --config ../config/proxy.yml
\end{bashcode}

\begin{notebox}[Go proxy status]
The Go proxy achieves feature parity with the Python proxy for all core pipeline
functions. Some auxiliary features (management API, analytics integration) are
still being finalized. See \code{docs/phases/PHASE\_15.md} for the current
parity status.
\end{notebox}

\section{Stopping the Stack}
\index{make stop}

To stop without removing data:
\begin{bashcode}
make stop
\end{bashcode}

To stop and remove all containers, networks, and volumes (clean slate):
\begin{bashcode}
make stop-clean
\end{bashcode}

\code{make stop} leaves Redis data and logs intact. \code{make stop-clean} removes
everything including Redis persistence.
